\section{Injecting water}
After removing atoms to create a fracture, and passivating the system, we are now ready to inject water into the fracture. To do this we use the technique of \emph{voxelation} (see \cref{sec:voxelation}). We first divide the system into voxels, find the voxels that make up the void, and then put water molecules in these voxels. The water density can then be controlled by the size of the voxels we use, and how many of the voxels we fill.

% \orangebox{
% \begin{itemize}
%     \item This could/should have been done using grand canonical ensemble / grand canonical monte carlo -- but this is computationally heavy?
%     \item The voxelation technique has problems in thin and/or skewed pores (not horizontal/vertical)
%     \item We want to have a water phase that is realistic (can be realistic with low or high density and pressure) and liquid
% \end{itemize}
% }

\subsection{Finding correct voxel size}
If we want to inject water with density $\rho$% [kg/m$^3$]
, we can find the voxel size we need from the molar mass of water, $M_\text{H$_2$O} = M = 0.0180158 \text{ kg/mol}$\todobo{source?}. We use the molar mass and wanted density to find the ``volume'' each water atom should occupy% , the unit used in the \hl{MD integrator/program and output files}
, as follows
\begin{align*}
    V 
    = \frac{ M\text{ [kg/mol]} \times \dfrac{1}{N_A \text{ [mol$^{-1}$]}}}{ \rho\text{ [kg/m$^3$]} } 
    = \frac{M}{\rho N_A}\text{[m$^3$]}
\end{align*}
where $N_A$ is the Avogadro constant. From here we find the size of the voxels by taking the cube root
\begin{align}
    l = \left(\frac{M}{\rho N_A} \right)^{1/3}\text{ m}.
    \label{eq:inject_water_voxel_size}
\end{align}
To get a water density approximately equal to $\rho$ we can then use a voxel size $l$, and put one water atom in each voxel. If we for example want to insert water with $\rho = 1000\text{ kg/m}^3$, approximately the density of water in room temperature\todobo{source?}, we get a voxel size of
\begin{align*}
    l = \left(\frac{0.0180158 \text{ kg/mol}}{1000\text{ kg/m$^3$} \times 6.0221 \times 10^{23}\text{ mol$^{-1}$}} \right)^{1/3} = 3.1 \text{ \AA}.
\end{align*}

\subsubsection{Voxel size in finite systems}
Since we have a finite system we usually can not use the exact voxel size we want, but we have to divide the system into an integer number of voxels. This means that we will not get the exact density we want if we fill all voxels. To remedy this we only fill the fraction of voxels to get the wanted density.

In practice we use the following procedure
\begin{itemize}
    \item Find the number of voxels to divide the system into (in each direction) via
    \begin{align*}
        n_i = \left\lceil \frac{L_i}{l_i} \right\rceil,
    \end{align*}
    where $L_i$ is the system size in dimension $i$ and $l_i$ is the voxel size calculated using \cref{eq:inject_water_voxel_size}. We use the ceiling-function to ensure that the actual voxel size we use is smaller than (or equal to) the voxel size we calculated. We find the actual voxel size via
    \begin{align*}
        \tilde l_i = \frac{L_i}{n_i}.
    \end{align*}
%     \item Mark all voxels within \hl{what radii?} as occupied.
    \item Divide the system into voxels of size $(\tilde l_x, \tilde l_y, \tilde l_z)$, and find the voxels that make up the void (see \cref{sec:inject_water_find_empty_voxels} for more on this).
    \item To find the ratio of voxels to put water molecules in we first calculate the density we would get if we filled all empty voxels using
    \begin{align*}
        \tilde\rho = \frac{M}{\tilde l_x \tilde l_y \tilde l_z N_A},
    \end{align*}
    and then find the number of voxels to fill as
    \begin{align*}
        \tilde N = N\frac{\tilde\rho}{\rho} = N\frac{\tilde l_x \tilde l_y \tilde l_z}{l_x l_y l_z},
    \end{align*}
    where $N$ is the total number of empty voxels.
    %\item \hl{Put a water molecule $\tilde N$ voxels}. 
    
    This can for example be done by looping over all voxels, drawing a random uniform number between 0 and 1 for each voxel, and putting a molecule in the voxel if the random number is smaller than $\tilde N / N$.
%     \begin{listing}[!htb]%
%         \begin{cppcode*}{gobble=12}
%             for (int i = 0; i < voxels.size(); i++)
%             {
%                 if (rand() < n_voxels_to_fill/n_voxels)
%                 {
%                     voxels[i].push_back(create_random_water_atom());
%                 }
%             }
%         \end{cppcode*}
%     \caption{%
%         \hl{FINISH THIS LISTING}.%
% %         \label{list:sampling}%
%     }%
%     \end{listing}%
\end{itemize}
The water molecules are inserted with the oxygen atom in the center of the voxel, and the two hydrogen atoms pointing in a random direction, but with an angle of $\sim 104.45^\circ$.

% \subsection[Finding empty/unoccupied voxels]{Finding \hl{empty/unoccupied} voxels\label{sec:inject_water_find_empty_voxels}}
\subsection{Identifying the voxels that make up the void\label{sec:inject_water_find_empty_voxels}}
The naive way of finding the voxels to fill with water is to just find which voxel each silicon and oxygen atom is in, and mark those as occupied. Using this method we found that we often got some empty voxels inside the silica matrix, which meant we got single water atoms trapped inside what was supposed to be the silica matrix. 

This can be explained if we compare the voxel size of $3.1 \text{ \AA}$ in water with a density of $\rho = 1000\text{ kg/m$^3$}$, as we found above, to the typical bond length between silica tetrahedra in solid silica, which is about \hl{???}. When we take into account the amorphous structure of silica we see that it is likely that we get some small pores with room for a water atom in between some of the silica tetrahedra.\todoao{better explanation}

To solve this we assign a radius to each atom type, and mark all voxels with the center of the voxel within this radius from an atom as occupied. The rest of the voxels should now be a good approximation to the void. See \cref{fig:inject_empty_voxel} for an illustration of this procedure. 
%
\begin{figure}[htpb]%
    \centering%
    \begin{subfigure}[b]{0.45\textwidth}%
        \includesvg[width=\textwidth, svgpath=./images/inject_water/]{drawing02}%
%         \caption{}%
        \caption{Marking only one voxel per atom as occupied.}%
%         \label{fig:pass_tet01}%
    \end{subfigure}%
    \hspace{0.05\textwidth}%
    \begin{subfigure}[b]{0.45\textwidth}%
        \includesvg[width=\textwidth, svgpath=./images/inject_water/]{drawing_radius02}%
%         \caption{}%
        \caption{Marking all voxels within radius from atom as occupied.}%
%         \label{fig:pass_tet02}%
    \end{subfigure}%
    \caption[
        To find voxels that make up the void/pore in we can either \textbf{a)} mark the voxel each existing atom belongs in as occupied, or \textbf{b)} mark all voxels within a radius from each atom as occupied. We can assign a different radius to each atom. We have illustrated using part of a silica tetrahedra, with one silicon atom (the large blue dot) and two oxygen atoms (the smaller red dot). The center of each voxel is marked by a dot 
    ]{%
        To find voxels that make up the void/pore in we can either \textbf{a)} mark the voxel each existing atom belongs in as occupied, or \textbf{b)} mark all voxels within a radius from each atom as occupied. We can assign a different radius to each atom. We have illustrated using part of a silica tetrahedra, with one silicon atom (\includegraphics[scale=0.8]{./images/inject_water/silicon.pdf}) and two oxygen atoms (\raisebox{0.3ex}{\includegraphics[scale=0.8]{./images/inject_water/oxygen.pdf}}). The center of each voxel is marked by a dot. %
        %\hl{Maybe change to v01 of images?} \hl{This voxel size is pretty tiny...}%
    }%
    \label{fig:inject_empty_voxel}%
\end{figure}%

\hl{A different solution to the problem of tiny pores inside the silica matrix is to remove all small clusters of voxels (where a ``small cluster of voxels'' would need to be defined), or maybe to use different voxel sizes for finding the void and filling the void with water.}\todoao{Fix or remove}

% \orangebox{
% \begin{itemize}
%     \item Remove tiny pores?
%     \item Give voxels random velocity based on wanted temperature? Now using zero!
%     \item Mark all voxels within distance from other atoms as occupied
%     \item Fill other voxels with H2O with random O-H orientation, but correct angle
%     \item Improvement: Use one voxel size in the beginning (to avoid one-voxel pores), and then use a smaller voxel size when injecting water
% \end{itemize}
% }

